\documentclass[reqno]{amsart}
\usepackage[T1]{fontenc}
\usepackage{amsmath,amssymb,amsthm}
\usepackage{xr-hyper}
\usepackage[hidelinks]{hyperref}
\externaldocument[sm-]{20261125T000000--slope-method__workbook}
\externaldocument[ec-]{20261124T223000--elliptic-curves-for-the-proof__workbook}
\externaldocument[st-]{20261124T214500--supersingular-theta-divisibility__workbook}
\externaldocument[cm-]{20261124T223100--complex-multiplication-and-inert-reduction__workbook}
\externaldocument[kl-]{20261125T009000--kronecker-limit-and-chowla-selberg__workbook}

% Chapter number; \ref{chap} gives it.
\newcounter{chap}
% Chapter numbers, generated from 20261125T012000--grh-workbook__guide.org
% by scripts/make-chapters.py.  Do not edit: change the order of the
% index and rerun the script (./preview-site.sh and CI do this).
%
% Usage in a sheet, after \newcounter{chap}:
%   % Chapter numbers, generated from 20261125T012000--grh-workbook__guide.org
% by scripts/make-chapters.py.  Do not edit: change the order of the
% index and rerun the script (./preview-site.sh and CI do this).
%
% Usage in a sheet, after \newcounter{chap}:
%   % Chapter numbers, generated from 20261125T012000--grh-workbook__guide.org
% by scripts/make-chapters.py.  Do not edit: change the order of the
% index and rerun the script (./preview-site.sh and CI do this).
%
% Usage in a sheet, after \newcounter{chap}:
%   \input{chapters}\chapnumber{<sheet ID>}
\makeatletter
\providecommand\chapters@set[2]{\@namedef{chapters@#1}{#2}}
\chapters@set{20261125T001000}{1}% zeros-to-ramp-sums
\chapters@set{20261125T001100}{2}% inert-primes-and-estimate-b
\chapters@set{20261125T005000}{3}% galois-fourier-and-averaging
\chapters@set{20261125T005500}{4}% columns-from-a-sphere
\chapters@set{20261125T000000}{5}% slope-method
\chapters@set{20261126T030100}{6}% a-twist-for-each-prime
\chapters@set{20261125T003000}{7}% the-chain-and-the-square-root
\chapters@set{20261125T010000}{8}% rank-lemma-and-zero-estimate
\chapters@set{20261125T010500}{9}% poincare-series-and-near-orthogonality
\chapters@set{20261124T223000}{10}% elliptic-curves-for-the-proof
\chapters@set{20261125T009500}{11}% fock-model-and-theta-intertwiners
\chapters@set{20261124T223100}{12}% complex-multiplication-and-inert-reduction
\chapters@set{20261125T008000}{13}% one-cm-determinant
\chapters@set{20261124T214500}{14}% supersingular-theta-divisibility
\chapters@set{20261124T214600}{15}% the-isotropic-plane
\chapters@set{20261125T013700}{16}% pairings-on-a-formal-group-of-height-two
\chapters@set{20261125T013500}{17}% hecke-nondegeneracy-on-the-tree
\chapters@set{20261126T030200}{18}% the-bound-is-attained
\chapters@set{20261125T009000}{19}% kronecker-limit-and-chowla-selberg
\chapters@set{20261125T002000}{20}% definite-theta-lifts
\chapters@set{20261125T015500}{21}% siegel-weil-in-the-smallest-case
\chapters@set{20261125T002100}{22}% definite-theta-lifts-and-norms
\chapters@set{20261125T016500}{23}% hecke-operators-and-theta-lifts
\chapters@set{20261125T016700}{24}% divisibility-along-the-chain
\chapters@set{20261125T004000}{25}% denominators-and-sha
\chapters@set{20261125T004100}{26}% tate-twists-and-sha
\chapters@set{20261126T020430}{27}% the-minus-part-by-gauss-sums
\chapters@set{20261125T014100}{28}% congruence-modules
\chapters@set{20261125T014500}{29}% eisenstein-congruences-and-selmer-classes
\chapters@set{20261125T014700}{30}% pseudocharacters-and-the-length-bound
\chapters@set{20261125T011000}{31}% galois-representations-black-box
\chapters@set{20261126T030300}{32}% the-data
\chapters@set{20261126T020210}{33}% bessel-poincare-series-and-their-gram-operators
\chapters@set{20261126T030400}{34}% the-seesaw-identity
\chapters@set{20261126T030500}{35}% divisibility-of-the-composite-lift
\chapters@set{20261126T030600}{36}% theorem-a-short-forms
\chapters@set{20261126T030700}{37}% the-determinant-of-a-chain
\chapters@set{20261126T030800}{38}% the-cost-of-one-lift
\chapters@set{20261126T030900}{39}% denominators-of-projections-to-lifts
\chapters@set{20261126T031000}{40}% twists-with-few-congruences
\chapters@set{20261126T031100}{41}% theorem-c-and-the-main-inequality
\providecommand\chapnumber[1]{%
  \@ifundefined{chapters@#1}%
    {\PackageError{chapters}{Sheet #1 is not listed in the index}%
      {List it in 20261125T012000--grh-workbook__guide.org\MessageBreak
       and run scripts/make-chapters.py.}}%
    {\edef\thechap{\@nameuse{chapters@#1}}}}
\makeatother
\chapnumber{<sheet ID>}
\makeatletter
\providecommand\chapters@set[2]{\@namedef{chapters@#1}{#2}}
\chapters@set{20261125T001000}{1}% zeros-to-ramp-sums
\chapters@set{20261125T001100}{2}% inert-primes-and-estimate-b
\chapters@set{20261125T005000}{3}% galois-fourier-and-averaging
\chapters@set{20261125T005500}{4}% columns-from-a-sphere
\chapters@set{20261125T000000}{5}% slope-method
\chapters@set{20261126T030100}{6}% a-twist-for-each-prime
\chapters@set{20261125T003000}{7}% the-chain-and-the-square-root
\chapters@set{20261125T010000}{8}% rank-lemma-and-zero-estimate
\chapters@set{20261125T010500}{9}% poincare-series-and-near-orthogonality
\chapters@set{20261124T223000}{10}% elliptic-curves-for-the-proof
\chapters@set{20261125T009500}{11}% fock-model-and-theta-intertwiners
\chapters@set{20261124T223100}{12}% complex-multiplication-and-inert-reduction
\chapters@set{20261125T008000}{13}% one-cm-determinant
\chapters@set{20261124T214500}{14}% supersingular-theta-divisibility
\chapters@set{20261124T214600}{15}% the-isotropic-plane
\chapters@set{20261125T013700}{16}% pairings-on-a-formal-group-of-height-two
\chapters@set{20261125T013500}{17}% hecke-nondegeneracy-on-the-tree
\chapters@set{20261126T030200}{18}% the-bound-is-attained
\chapters@set{20261125T009000}{19}% kronecker-limit-and-chowla-selberg
\chapters@set{20261125T002000}{20}% definite-theta-lifts
\chapters@set{20261125T015500}{21}% siegel-weil-in-the-smallest-case
\chapters@set{20261125T002100}{22}% definite-theta-lifts-and-norms
\chapters@set{20261125T016500}{23}% hecke-operators-and-theta-lifts
\chapters@set{20261125T016700}{24}% divisibility-along-the-chain
\chapters@set{20261125T004000}{25}% denominators-and-sha
\chapters@set{20261125T004100}{26}% tate-twists-and-sha
\chapters@set{20261126T020430}{27}% the-minus-part-by-gauss-sums
\chapters@set{20261125T014100}{28}% congruence-modules
\chapters@set{20261125T014500}{29}% eisenstein-congruences-and-selmer-classes
\chapters@set{20261125T014700}{30}% pseudocharacters-and-the-length-bound
\chapters@set{20261125T011000}{31}% galois-representations-black-box
\chapters@set{20261126T030300}{32}% the-data
\chapters@set{20261126T020210}{33}% bessel-poincare-series-and-their-gram-operators
\chapters@set{20261126T030400}{34}% the-seesaw-identity
\chapters@set{20261126T030500}{35}% divisibility-of-the-composite-lift
\chapters@set{20261126T030600}{36}% theorem-a-short-forms
\chapters@set{20261126T030700}{37}% the-determinant-of-a-chain
\chapters@set{20261126T030800}{38}% the-cost-of-one-lift
\chapters@set{20261126T030900}{39}% denominators-of-projections-to-lifts
\chapters@set{20261126T031000}{40}% twists-with-few-congruences
\chapters@set{20261126T031100}{41}% theorem-c-and-the-main-inequality
\providecommand\chapnumber[1]{%
  \@ifundefined{chapters@#1}%
    {\PackageError{chapters}{Sheet #1 is not listed in the index}%
      {List it in 20261125T012000--grh-workbook__guide.org\MessageBreak
       and run scripts/make-chapters.py.}}%
    {\edef\thechap{\@nameuse{chapters@#1}}}}
\makeatother
\chapnumber{<sheet ID>}
\makeatletter
\providecommand\chapters@set[2]{\@namedef{chapters@#1}{#2}}
\chapters@set{20261125T001000}{1}% zeros-to-ramp-sums
\chapters@set{20261125T001100}{2}% inert-primes-and-estimate-b
\chapters@set{20261125T005000}{3}% galois-fourier-and-averaging
\chapters@set{20261125T005500}{4}% columns-from-a-sphere
\chapters@set{20261125T000000}{5}% slope-method
\chapters@set{20261126T030100}{6}% a-twist-for-each-prime
\chapters@set{20261125T003000}{7}% the-chain-and-the-square-root
\chapters@set{20261125T010000}{8}% rank-lemma-and-zero-estimate
\chapters@set{20261125T010500}{9}% poincare-series-and-near-orthogonality
\chapters@set{20261124T223000}{10}% elliptic-curves-for-the-proof
\chapters@set{20261125T009500}{11}% fock-model-and-theta-intertwiners
\chapters@set{20261124T223100}{12}% complex-multiplication-and-inert-reduction
\chapters@set{20261125T008000}{13}% one-cm-determinant
\chapters@set{20261124T214500}{14}% supersingular-theta-divisibility
\chapters@set{20261124T214600}{15}% the-isotropic-plane
\chapters@set{20261125T013700}{16}% pairings-on-a-formal-group-of-height-two
\chapters@set{20261125T013500}{17}% hecke-nondegeneracy-on-the-tree
\chapters@set{20261126T030200}{18}% the-bound-is-attained
\chapters@set{20261125T009000}{19}% kronecker-limit-and-chowla-selberg
\chapters@set{20261125T002000}{20}% definite-theta-lifts
\chapters@set{20261125T015500}{21}% siegel-weil-in-the-smallest-case
\chapters@set{20261125T002100}{22}% definite-theta-lifts-and-norms
\chapters@set{20261125T016500}{23}% hecke-operators-and-theta-lifts
\chapters@set{20261125T016700}{24}% divisibility-along-the-chain
\chapters@set{20261125T004000}{25}% denominators-and-sha
\chapters@set{20261125T004100}{26}% tate-twists-and-sha
\chapters@set{20261126T020430}{27}% the-minus-part-by-gauss-sums
\chapters@set{20261125T014100}{28}% congruence-modules
\chapters@set{20261125T014500}{29}% eisenstein-congruences-and-selmer-classes
\chapters@set{20261125T014700}{30}% pseudocharacters-and-the-length-bound
\chapters@set{20261125T011000}{31}% galois-representations-black-box
\chapters@set{20261126T030300}{32}% the-data
\chapters@set{20261126T020210}{33}% bessel-poincare-series-and-their-gram-operators
\chapters@set{20261126T030400}{34}% the-seesaw-identity
\chapters@set{20261126T030500}{35}% divisibility-of-the-composite-lift
\chapters@set{20261126T030600}{36}% theorem-a-short-forms
\chapters@set{20261126T030700}{37}% the-determinant-of-a-chain
\chapters@set{20261126T030800}{38}% the-cost-of-one-lift
\chapters@set{20261126T030900}{39}% denominators-of-projections-to-lifts
\chapters@set{20261126T031000}{40}% twists-with-few-congruences
\chapters@set{20261126T031100}{41}% theorem-c-and-the-main-inequality
\providecommand\chapnumber[1]{%
  \@ifundefined{chapters@#1}%
    {\PackageError{chapters}{Sheet #1 is not listed in the index}%
      {List it in 20261125T012000--grh-workbook__guide.org\MessageBreak
       and run scripts/make-chapters.py.}}%
    {\edef\thechap{\@nameuse{chapters@#1}}}}
\makeatother
\chapnumber{20261125T008000}
\renewcommand{\thesection}{\thechap.\arabic{section}}

\theoremstyle{definition}
\newtheorem{problem}{Problem}[section]
\theoremstyle{plain}
\newtheorem{theorem}{Theorem}[section]

\newcommand{\F}{\mathbb{F}}
\newcommand{\Z}{\mathbb{Z}}
\newcommand{\Q}{\mathbb{Q}}
\newcommand{\C}{\mathbb{C}}
\newcommand{\Zi}{\Z[i]}

\setlength{\emergencystretch}{2em}

\title{One CM determinant: Mertens modulo 4}
\author{}
\date{November 25, 2026}

\begin{document}
\maketitle
\refstepcounter{chap}\label{chap}

\paragraph{Prerequisites.}
Legendre's formula \(v_p(r!) = \sum_{e \ge 1} \lfloor r/p^e \rfloor\), Stirling's formula, Chebyshev's bound (Problem~\ref{sm-pr:binomial} in Chapter~\ref{sm-chap}), and the arithmetic of \(\Zi\).  The invariant differential and complex uniformization of an elliptic curve \cite[\S III.5, Ch.~VI]{Si1}, and supersingular reduction \cite[\S V.4]{Si1}.  The Jacobi theta function \(\theta_1\) \cite[\S20.2]{DLMF}, Cauchy's inequality and the argument principle.  Section~\ref{sec:linear} also uses the prime number theorem for arithmetic progressions \cite[\S11.3]{MV}.  From Chapter~\ref{ec-chap}, \emph{Complex tori, theta functions and periods}: Jacobi's function \(\vartheta\) (Section~\ref{ec-sec:theta}).  From Chapter~\ref{cm-chap}, \emph{Complex multiplication and reduction at an inert prime}: the congruence \(D^p \equiv \lambda D \pmod p\) for a formal group (Section~\ref{cm-sec:formal}) and the supersingular reduction of \(y^2 = x^3 - x\) at the primes inert in \(\Zi\) (Section~\ref{cm-sec:reduction}).  Chapter~\ref{st-chap}, \emph{The invariant derivation at a supersingular prime}, treats the derivation \(D^p\) for the lemniscate.

\section{Mertens from the factorial}\label{sec:factorial}

Let \(\theta(X) := \sum_{p \le X} \log p\).  How much of \(\log r!\) comes from each prime?  Chebyshev's bound \(\theta(2n) - \theta(n) \le n \log 4\) is Problem~\ref{sm-pr:binomial} of Chapter~\ref{sm-chap}.

\begin{problem}\label{pr:chebyshev}
\(\theta(X) < 2X\log 4\) for \(X \ge 1\).
\end{problem}

\begin{problem}\label{pr:factorial}
For \(r \ge 1\),
\(-\theta(r) \le \log r! - r \sum_{p \le r} \frac{\log p}{p} \le r \sum_p \frac{\log p}{p(p-1)}\).
\end{problem}

With Stirling's formula, Problems~\ref{pr:chebyshev} and~\ref{pr:factorial} prove the following \cite{Mer}.

\begin{theorem}[Mertens]\label{thm:mertens}
\(\sum_{p \le X} \log p/p = \log X + O(1)\) for \(X \ge 2\).
\end{theorem}

How much of \(\log r!\) comes from each odd class modulo~\(4\)?  Let
\[
  S_a(X) := \sum_{p \le X,\ p \equiv a \ (4)} \frac{\log p}{p}.
\]

\begin{theorem}\label{thm:mod4}
\(S_a(X) = \frac12 \log X + O(1)\) for \(a = 1, 3\) and \(X \ge 2\).
\end{theorem}

This is classical, and the usual proof rests on \(L(1, \chi_{-4}) \ne 0\) \cite[Ch.~4]{MV}.  Since \(S_1 + S_3 = \log X + O(1)\), two upper bounds suffice.  We replace the factorial by one determinant for each class and take \(r = \lfloor X \rfloor\).

\section{The Gaussian Vandermonde}\label{sec:vandermonde}

Let \(z_1, \dots, z_r\) be distinct Gaussian integers with \(0 \le \operatorname{Re} z, \operatorname{Im} z < \lceil \sqrt r \rceil\), and \(V := \prod_{j < k} (z_k - z_j)\).

\begin{problem}\label{pr:V-size}
\(\lvert V \rvert^2\) is a positive integer, and \(\log \lvert V \rvert^2 \le \binom r2 \log(2r)\).
\end{problem}

\begin{problem}\label{pr:V-split}
If \(p \equiv 1 \pmod 4\), then \(v_p(\lvert V \rvert^2) \ge r^2/p - r\).
\end{problem}

Summing Problem~\ref{pr:V-split} with weight \(\log p\) over \(p \le r\) gives \(r^2 S_1(r) \le \log \lvert V \rvert^2 + r\theta(r)\).  With Problems~\ref{pr:chebyshev} and~\ref{pr:V-size} this proves the following.  The counting of coincident pairs follows \cite[\S2]{PV}.

\begin{theorem}\label{thm:split}
\(S_1(X) \le \frac12 \log X + O(1)\).
\end{theorem}

\begin{problem}\label{pr:V-inert}
If \(p \equiv 3 \pmod 4\), then \(v_p(\lvert V \rvert^2) \ge r^2/p^2 - r\).
\end{problem}

The weights \(\log p/p^2\) have a convergent sum, so Problem~\ref{pr:V-inert} says nothing about \(S_3\).  An inert prime has residue field \(\F_{p^2}\), and Gaussian integers coincide modulo \(p\) only with frequency \(1/p^2\).  To see an inert prime with frequency \(1/p\) we use the curve with complex multiplication by \(\Zi\), which is supersingular exactly at these primes.  Its invariant derivation \(D\) satisfies \(D^p \equiv 0 \pmod p\), so every block of \(p\) derivatives gains a factor~\(p\), as every block of \(p\) consecutive integers does in~\(r!\).

\section{The derivation at an inert prime}\label{sec:derivation}

Let \(E \colon y^2 = x^3 - x\), \(A := \Z[x, y]/(y^2 - x^3 + x)\), \(\omega := dx/(2y)\), and \(P := (0, 0)\).  The derivation \(D\) of \(A\) dual to \(\omega\) is
\[
  Dx = 2y, \qquad Dy = 3x^2 - 1.
\]
The automorphism \([i](x, y) := (-x, iy)\) fixes \(P\) and pulls \(\omega\) back to \(i\omega\).  Let \(d_k := (D^k x)(P)\).

\begin{problem}\label{pr:dk}
The \(d_k\) are integers, and \(d_k = 0\) unless \(k \equiv 2 \pmod 4\).  Moreover \(d_2 = -2\), \(d_6 = 144\), \(d_{10} = -96768\).
\end{problem}

Let \(s = 1 + \frac14 t^4 + \dotsb\) be the root of \(s^4 = 1 + t^4\) with constant term~\(1\).  For odd \(p\) it lies in \(\Z_{(p)}[[t]]\).

\begin{problem}\label{pr:fermat}
The point \(x = -t^2/s^2\), \(y = -t/s^3\) lies on \(E\), \(t\) is a parameter at \(P\), and
\[
  \omega = \frac{dt}{s^3} = \sum_{h \ge 0} (-1)^h b_h \, t^{4h} \, dt,
  \qquad
  b_h := \frac{\prod_{j=1}^h (4j - 1)}{4^h h!}.
\]
\end{problem}

Thus \(E\) is a quotient of the Fermat quartic \(s^4 = 1 + t^4\), and its differential is a series in \(t^4\).

\begin{problem}\label{pr:Dp-cD}
Let \(p\) be an odd prime.  On \(A/pA\) the operator \(D^p\) is a derivation, and \(D^p = cD\) for a constant \(c \in \F_p\).
\end{problem}

On the formal group of \(E\) this is the congruence \(D^p \equiv \lambda D \pmod p\) of Section~\ref{cm-sec:formal}, with \(\lambda \equiv c\).

\begin{problem}\label{pr:c-zero}
If \(p \equiv 3 \pmod 4\), then \(c = 0\).
\end{problem}

This is the vanishing of the Hasse invariant at a supersingular prime \cite[Thm.~V.4.1(a)]{Si1}.

\begin{problem}\label{pr:dk-div}
If \(p \equiv 3 \pmod 4\), then \(p^{\lfloor k/p \rfloor}\) divides \((D^k f)(P)\) for all \(f \in A\) and \(k \ge 0\).  For instance \(3^2 \mid d_6\) and \(7^2 \mid d_{14} = 268240896\).
\end{problem}

\section{The determinant of derivatives}\label{sec:determinant}

Write \(E(\C) = \C/\Lambda\) with \(dz = \omega\) \cite[Ch.~VI]{Si1}, and let \(w(z)\) be the \(x\)-coordinate of \(P + z\).  Then \(w(z) = \wp(z + u_0)\) for a lift \(u_0 \notin \Lambda\) of \(P\), \(D = d/dz\), \(d_k = w^{(k)}(0)\), and
\[
  (w')^2 = 4(w^3 - w), \qquad w(0) = w'(0) = 0, \qquad w''(0) = -2.
\]
For \(r \ge 1\), let \(M\) be the matrix whose row of index \(k \ge 0\) lists the \(k\)th derivatives at \(0\) of
\begin{equation}\label{eq:columns}
  \frac{z^j}{j!} \quad (0 \le j < r), \qquad \frac{z^j}{j!}\, w(z) \quad (0 \le j < r).
\end{equation}
The first \(r\) columns are the identity followed by zeros.

\begin{problem}\label{pr:entries}
The entry of \(M\) in row \(k\) and column \(z^j w/j!\) is \(\binom kj d_{k - j}\) for \(k \ge j\), and \(0\) otherwise.
\end{problem}

A relation among the columns, restricted to rows \(0, \dots, 4r\), is a pair of polynomials \(a, b\) of degree \(< r\), not both zero, with \(a + bw = O(z^{4r+1})\).

\begin{problem}\label{pr:Q}
Let \(Q(a, b) := (a'b - ab')^2 + 4ab(a^2 - b^2)\).  Then \(Q(a, b) = b^4\bigl((h')^2 - 4(h^3 - h)\bigr)\) with \(h = -a/b\), when \(b \ne 0\).  If \(a + bw = O(z^{4r+1})\) with \(\deg a, \deg b < r\), then \(Q(a, b) = 0\).
\end{problem}

\begin{problem}\label{pr:no-rational}
No nonconstant rational function \(h \in \C(z)\) satisfies \((h')^2 = 4(h^3 - h)\).
\end{problem}

\begin{problem}\label{pr:constant}
If \(h \in \{0, \pm1\}\) and \(b \ne 0\) has degree \(< r\), then \(b(w - h)\) vanishes at \(0\) to order at most \(r + 1\).
\end{problem}

Problems~\ref{pr:Q}--\ref{pr:constant} prove the following.

\begin{theorem}\label{thm:rank}
The rows of \(M\) of index \(0, \dots, 4r\) have rank \(2r\).
\end{theorem}

In order of increasing \(k\), keep each row of \(M\) that is independent of its predecessors, until \(2r\) rows are kept.  Their indices are \(I = \{0, \dots, r-1, k_1, \dots, k_r\}\) with \(r \le k_1 < \dots < k_r \le 4r\).  Expanding along the identity block,
\[
  N_r := \det M_I = \det\Bigl[\tbinom{k_i}{j} d_{k_i - j}\Bigr]_{1 \le i \le r,\ 0 \le j < r}
\]
is a nonzero integer.  Let \(L := \sum_i k_i - \binom r2\), the total order \(\sum_i (k_i - j_i)\) of the derivatives in each term of this determinant.

\begin{problem}\label{pr:Nr-small}
For \(r = 1, 2, 3\) the new rows are \(\{2\}\), \(\{2, 3\}\), \(\{3, 4, 6\}\), and \(N_r = -2\), \(12\), \(2^7 3^4\).
\end{problem}

\begin{problem}\label{pr:Nr-div}
\(r^2 \le L \le 4r^2\), and \(v_p(N_r) \ge L/p - r\) for \(p \equiv 3 \pmod 4\).
\end{problem}

\section{The size of the determinant}\label{sec:size}

The entries \(d_k\) grow like \(k!\) times a power of the distance from \(0\) to the nearest pole of \(w\).  To bound \(N_r\) we first remove the poles.

\begin{problem}\label{pr:triangular}
Let \(H\) be holomorphic near \(0\) with \(H(0) = 1\).  Multiplying all the functions \eqref{eq:columns} by \(H\) does not change \(\det M_I = N_r\).
\end{problem}

Write \(\Lambda = \lambda(\Z + \tau\Z)\) with \(\operatorname{Im}\tau > 0\), and let
\[
  \theta_\tau(v) := 2\sum_{n \ge 0} (-1)^n e^{\pi i\tau(n + 1/2)^2} \sin\bigl((2n + 1)\pi v\bigr),
\]
so that \(\theta_\tau(v) = \theta_1(\pi v \mid \tau)\) \cite[(20.2.1)]{DLMF}.  In terms of Jacobi's function \(\vartheta\) of Section~\ref{ec-sec:theta}, \(\theta_\tau(v) = -i\, e^{\pi i\tau/4 + \pi i v}\, \vartheta\bigl(v + \tfrac{1+\tau}{2}\bigr)\).  It is odd, \(\theta_\tau(v + 1) = -\theta_\tau(v)\), and \(\theta_\tau(v + \tau) = -e^{-\pi i\tau - 2\pi i v}\theta_\tau(v)\).

\begin{problem}\label{pr:theta-zeros}
The zeros of \(\theta_\tau\) are the points of \(\Z + \tau\Z\), and they are simple.
\end{problem}

Let \(H(z) := \theta_\tau\bigl((z + u_0)/\lambda\bigr)^2 / \theta_\tau(u_0/\lambda)^2\).

\begin{problem}\label{pr:H-growth}
\(H\) and \(Hw\) are entire, and \(\lvert H(z) \rvert, \lvert H(z) w(z) \rvert \le e^{C(1 + \lvert z \rvert^2)}\) with \(C\) independent of \(z\).
\end{problem}

\begin{problem}\label{pr:Nr-size}
For every \(R > 0\), \(\lvert N_r \rvert \le (2r)!\,(4r/R)^L e^{2Cr(1 + R^2)}\).  Hence \(\log \lvert N_r \rvert \le \frac12 L \log r + O(r^2)\).
\end{problem}

Summing Problem~\ref{pr:Nr-div} with weight \(\log p\) over \(p \le r\) gives \(L S_3(r) \le \log\lvert N_r \rvert + r\theta(r)\).  By Problems~\ref{pr:chebyshev} and~\ref{pr:Nr-size}, and \(L \ge r^2\), this proves the following.

\begin{theorem}\label{thm:inert}
\(S_3(X) \le \frac12 \log X + O(1)\).
\end{theorem}

Theorems~\ref{thm:mertens}, \ref{thm:split} and~\ref{thm:inert} prove Theorem~\ref{thm:mod4}.  The proof never uses \(L(1, \chi_{-4}) \ne 0\).

\section{The linear term}\label{sec:linear}

Both determinants lose \(O(r^2)\): the floors \(\lfloor k/p \rfloor\) cost \(r\theta(r)\), and the bounds for \(\lvert V \rvert\) and \(\lvert N_r \rvert\) are crude at order \(r^2\).  So they give \(S_a\) up to \(O(1)\) and say nothing of the constant term.  For \(\chi = \chi_{-4}\) that term involves
\[
  c_\chi := \frac{L'}{L}(1, \chi) = \gamma + \log\frac{\pi^2}{2\varpi^2} = 0.2456\ldots,
  \qquad
  \varpi := 2\int_0^1 \frac{dt}{\sqrt{1 - t^4}},
\]
where the second expression follows from Kronecker's limit formula and the value of \(\varpi\) in terms of \(\Gamma(\frac14)\) (Theorem~\ref{ec-thm:LL}); it is the case \(\Q(i)\) of the formula of Chowla and Selberg \cite{CS}.  How close can a determinant come to \(c_\chi\)?

On the Vandermonde side the floors can be kept exactly.  For a prime power \(q\) let \(w_q(k) := \sum_{e \ge 1} \lfloor k/q^e \rfloor\), and let \(k!_{\Zi} := \prod_{\mathfrak p} \mathfrak p^{w_{N\mathfrak p}(k)}\), the \(k\)th factorial of \(\Zi\) in the sense of Pólya and Bhargava.  For any \(n\) Gaussian integers the Vandermonde product is divisible by \(\prod_{k < n} k!_{\Zi}\) \cite{Bh}.  Write \(\lvert k!_{\Zi} \rvert\) for the absolute value of a generator.

\begin{problem}\label{pr:bhargava}
\(\log \dfrac{k!}{\lvert k!_{\Zi} \rvert} = \frac12 w_2(k) \log 2 + \sum_{p \equiv 3 \ (4)} \log p \sum_{e \text{ odd}} \Bigl\lfloor \dfrac{k}{p^e} \Bigr\rfloor\).  For \(k = 6\) both sides are \(\log 36\).
\end{problem}

Let \(\beta(n) := \sum_{k < n} \log\lvert k!_{\Zi} \rvert\).

\begin{problem}\label{pr:beta-sum}
\(\beta(n) = \frac12 \sum_{N < n} (n - N) \bigl( \log N + \sum_{p^e \mid N} \chi(p^e) \log p \bigr)\), the inner sum over prime powers \(p^e > 1\).
\end{problem}

The prime number theorem for \(\chi\) \cite[\S11.3]{MV} gives \(\sum_{p^e \le x} \chi(p^e) \log p = o(x)\) and \(\sum_{p^e \le x} \chi(p^e) \log p / p^e = -c_\chi + o(1)\).

\begin{problem}\label{pr:beta}
\(\beta(n) = \frac{n^2}4 \log n - \frac{n^2}{4}\bigl(\frac32 + c_\chi\bigr) + o(n^2)\).
\end{problem}

We take the points in a disc.

\begin{problem}\label{pr:disc}
Let \(\mu\) be the uniform probability measure on the disc \(\lvert z \rvert \le R\).  For \(\lvert z \rvert \le R\), \(\int \log\lvert z - u \rvert \, d\mu(u) = \log R - (R^2 - \lvert z \rvert^2)/(2R^2)\), and hence \(\iint \log\lvert z - u \rvert \, d\mu \, d\mu = \log R - \frac14\).
\end{problem}

\begin{problem}\label{pr:Vn}
Let \(V_n\) be the Vandermonde product of the \(n\) Gaussian integers nearest \(0\).  Then \(\log\lvert V_n \rvert = \frac{n^2}{4} \log\frac n\pi - \frac{n^2}8 + o(n^2)\).
\end{problem}

\begin{problem}\label{pr:ell-lower}
\(\beta(n) \le \log\lvert V_n \rvert\) implies \(c_\chi \ge \log\pi - 1 = 0.1447\ldots\).
\end{problem}

So even with exact divisibility the Vandermonde loses \(c_\chi - \log\pi + 1 = 0.1008\ldots\) in the linear term.  On the side of derivatives one may build Padé-type determinants from the Taylor coefficients of the lemniscatic sine, whose poles are the points of \(\varpi\bigl(\frac{1+i}2 + \Zi\bigr)\).  We expect, from computation, that they give \(c_\chi \le \log(4\pi/\varpi^2) = 0.6031\ldots\), short by the absolute constant \(\log(8/\pi) - \gamma = 0.3574\ldots\); we do not know a proof of the size estimate this needs.  In both directions the loss comes from a disc of Gaussian integers, as nodes on one side and as poles on the other.  An error \(O(1)\) in \(S_a\) is Chebyshev's strength.  We expect that one determinant is not enough to control the constant term, let alone the error terms that the Riemann hypothesis requires.

In Alp\"oge's proof the constant term is not estimated; it cancels.  The size of a derivative of order \(a\) of a theta function with complex multiplication by \(K = \Q(\sqrt{-d})\) is measured against a period whose logarithm contains \(-c_1\), where \(c_1 = -1 - \gamma + \operatorname{Re} L'(1, \chi_K)/L(1, \chi_K)\) is also the constant in the expected value \(a(\log(a/\bar\kappa) + c_1)\) of the sum over primes that the divisibility produces; here \(\bar\kappa\) is a geometric mean of orders of derivatives.  The difference of the two amounts therefore contains no \(c_1\).  Chapter~\ref{kl-chap} explains why a period of an elliptic curve with complex multiplication involves \(L'(1, \chi_K)/L(1, \chi_K)\).

\paragraph{In the Lean formalization.}
In the rewritten Lean proof \cite{GL} this is \texttt{jetSize\_sub\_mainPrimeSum\_le}, in \texttt{New/UpperBound/Size/ArchimedeanSize.lean} of the readable edition: the size \(\log a! + a\log(M/\Omega)\) of a derivative of order \(a\), less the main amount \(a(\log(a/\bar\kappa) + c_1)\), is at most \(a\log(\bar\kappa M\,\mathrm{rd}_E/2\pi) + 1 + \frac12\log a\), where \(\Omega = e^{\gamma - \operatorname{Re} L'(1,\chi_K)/L(1,\chi_K)} \cdot 2\pi/\mathrm{rd}_E\), \(\mathrm{rd}_E\) is the root discriminant of a CM field \(E \supset K\), and \(M\) is the product of the two primes of the level.

\subsection*{Exercises}

\begin{enumerate}
\item If \(p \equiv 1 \pmod 4\) and \(h = (p-1)/4\), then \(c \equiv (-1)^h b_h \not\equiv 0 \pmod p\) in Problem~\ref{pr:Dp-cD}.  This is the coefficient of \(x^{p-1}\) in \((x^3 - x)^{(p-1)/2}\): \(c = 3, 6, 2\) for \(p = 5, 13, 17\).
\item \(N_4 = 2^{12}\, 3^5 \cdot 5 \cdot 7\), with new rows \(4, 5, 6, 7\).
\end{enumerate}

\section*{Hints}

\subsection*{\S\ref{sec:factorial}}
\begin{description}
\item[\ref{pr:chebyshev}] Sum Problem~\ref{sm-pr:binomial} over \(n\) near \(X/2^m\).
\item[\ref{pr:factorial}] Expand \(\log r!\) by Legendre's formula; compare \(\lfloor r/p \rfloor\) with \(r/p\), and bound the terms with \(e \ge 2\) by a geometric series.
\end{description}

\subsection*{\S\ref{sec:vandermonde}}
\begin{description}
\item[\ref{pr:V-size}] \(\lvert V \rvert^2\) is the norm of \(V\).
\item[\ref{pr:V-split}] Let \((p) = \mathfrak p\bar{\mathfrak p}\), and count the pairs \(j < k\) with \(z_j \equiv z_k\) modulo \(\mathfrak p\) and modulo \(\bar{\mathfrak p}\).
\item[\ref{pr:V-inert}] Now \(\Zi/p \cong \F_{p^2}\), and \(p \mid z\) gives \(p^2 \mid \lvert z \rvert^2\).
\end{description}

\subsection*{\S\ref{sec:derivation}}
\begin{description}
\item[\ref{pr:dk}] Compute \(D^k x\) in \(A\), using \(y^2 = x^3 - x\).  For the vanishing, compare \(D\) with its transform under the automorphism \([i]\), which fixes \(P\).
\item[\ref{pr:fermat}] \(4s^3 s' = 4t^3\).  Then \(dx = -2t\, dt/s^6\).  Expand \(s^{-3} = (1 + t^4)^{-3/4}\) binomially; \(y = -t + O(t^5)\).
\item[\ref{pr:Dp-cD}] Leibniz's rule, and \(p \mid \binom pi\) for \(0 < i < p\).  The derivation \(D\) is dual to the regular nowhere vanishing \(\omega\), so \(D^p = cD\) with \(c\) a function regular on all of \(E\).
\item[\ref{pr:c-zero}] \(c \equiv \lambda \equiv c_p = A_p\) by Problems~\ref{cm-pr:lambda} and~\ref{cm-pr:log}; then Problem~\ref{cm-pr:hasse-x3x}.
\item[\ref{pr:dk-div}] Lift \(D^p A \subseteq pA\) and iterate.
\end{description}

\subsection*{\S\ref{sec:determinant}}
\begin{description}
\item[\ref{pr:Q}] For the second part, substitute \(a = -bw + O(z^{4r+1})\), and compare the order of vanishing at \(0\) with the degree.
\item[\ref{pr:no-rational}] The map \(z \mapsto (h, h'/2)\) would extend to \(\mathbb P^1 \to E\) and pull \(\omega\) back to \(dz\).  Look at \(z = \infty\).
\item[\ref{pr:constant}] \(w\) has a double zero at \(0\).
\item[\ref{pr:Nr-small}] Use Problem~\ref{pr:entries} and the values in Problem~\ref{pr:dk}.
\item[\ref{pr:Nr-div}] Use \(k_i \ge r + i - 1\).  Apply Problem~\ref{pr:dk-div} to each term of the determinant, and use \(\lfloor m/p \rfloor > m/p - 1\).
\end{description}

\subsection*{\S\ref{sec:size}}
\begin{description}
\item[\ref{pr:triangular}] Leibniz's rule, and the way the rows \(I\) were chosen.
\item[\ref{pr:theta-zeros}] Integrate \(\theta_\tau'/\theta_\tau\) around a period parallelogram; the sides cancel up to the jump \(-2\pi i\) under \(v \mapsto v + \tau\).  Use oddness for the location.
\item[\ref{pr:H-growth}] The double zeros of \(H\) sit at the double poles of \(w\).  Bound \(H\), \(Hw\) on one parallelogram, and translate using the functional equation of \(\theta_\tau\).
\item[\ref{pr:Nr-size}] Apply Problem~\ref{pr:triangular} and Cauchy's inequality to each entry, then choose \(R\).
\end{description}

\subsection*{\S\ref{sec:linear}}
\begin{description}
\item[\ref{pr:bhargava}] \(2 = -i(1 + i)^2\); a split \(p\) has two primes of norm \(p\), and an inert \(p\) is prime of norm \(p^2\).
\item[\ref{pr:beta-sum}] By Problem~\ref{pr:bhargava}, \(\log\lvert k!_{\Zi} \rvert = \sum_{p^e} \frac{1 + \chi(p^e)}2 \log p \, \lfloor k/p^e \rfloor\).  Then \(\sum_{k < n} \lfloor k/d \rfloor = \sum_{N < n,\, d \mid N} (n - N)\).
\item[\ref{pr:beta}] \(\sum_{1 \le m < x} (x - m) = \frac{x^2}2 - \frac x2 + \frac12 \{x\}(1 - \{x\})\).  For the last term with \(x = n/d\), treat \(d \le \delta n\) trivially, and use partial summation on each interval where \(\lfloor n/d \rfloor\) is constant.
\item[\ref{pr:disc}] \(\frac1{2\pi}\int_0^{2\pi} \log\lvert z - \rho e^{i\phi} \rvert \, d\phi = \log\max(\lvert z \rvert, \rho)\).
\item[\ref{pr:Vn}] Scale by \(R = \sqrt{n/\pi}\) and compare the sum over pairs with the integral of Problem~\ref{pr:disc}; compare \(\sum \log(\lvert z - u \rvert/R)\) with the integral; pairs at distance \(< \varepsilon R\) contribute \(O(\varepsilon^2 \log(1/\varepsilon)\, n^2)\).
\item[\ref{pr:ell-lower}] Compare the coefficients of \(n^2\).
\end{description}

\begin{thebibliography}{DLMF}

\bibitem[Bh]{Bh}
M.~Bhargava, \emph{\(P\)-orderings and polynomial functions on arbitrary subsets of Dedekind rings}, J. reine angew. Math. \textbf{490} (1997), 101--127.

\bibitem[CS]{CS}
A.~Selberg and S.~Chowla, \emph{On Epstein's zeta-function}, J. reine angew. Math. \textbf{227} (1967), 86--110.

\bibitem[DLMF]{DLMF}
\emph{NIST Digital Library of Mathematical Functions}, \url{https://dlmf.nist.gov/}, Chapter~20.

\bibitem[GL]{GL} L.~Sernau, \emph{GRH for Dirichlet \(L\)-functions, in Lean~4}, a shorter rewriting of the Lean proof, text by Claude, \url{https://github.com/lsernau/grh-lean}, commit \texttt{72fead1}, readable edition \texttt{grh/}.
\bibitem[Mer]{Mer}
F.~Mertens, \emph{Ein Beitrag zur analytischen Zahlentheorie}, J. reine angew. Math. \textbf{78} (1874), 46--62.

\bibitem[MV]{MV}
H.~L.~Montgomery and R.~C.~Vaughan, \emph{Multiplicative Number Theory I. Classical Theory}, Cambridge Stud. Adv. Math. 97, Cambridge University Press, 2007.

\bibitem[PV]{PV}
V.~V.~Volkov and F.~V.~Petrov, \emph{On the interpolation of integer-valued polynomials}, J. Number Theory \textbf{133} (2013), 4224--4232; arXiv:1108.3212.

\bibitem[Si1]{Si1}
J.~H.~Silverman, \emph{The Arithmetic of Elliptic Curves}, 2nd ed., Grad. Texts in Math. 106, Springer, 2009.

\end{thebibliography}

\end{document}
